% STATUT : À RÉDIGER, résultats déjà produits (euro_cascade_model.py,
% scripts SCR multi-états, distribution, Delta_DORA).
\chapter{De la cascade au capital}
\label{ch:severite}

\textit{[Chapitre à rédiger.]}

% PLAN :
%  1. Chaîne complète : fréquence d'entrée -> pilier d'amorce -> cascade
%     (ensemble de piliers touchés) -> sévérité par pilier -> agrégation
%     -> distribution de perte annuelle -> SCR à 99,5 %.
%  2. Sévérité en euros : GPD au-delà du seuil, corps empirique en deçà,
%     plafond de réassurance. Justifier chaque choix.
%  3. État de conformité par pilier (NC / PC / C) et sa dynamique de
%     Markov (validé avec Hugo, remplace la latente statique). Le gain de
%     propagation g et la matérialité dépendent de l'état.
%  4. Résultats : SCR_DORA en DISTRIBUTION (jamais un point), Delta_DORA
%     entre profil conforme et non conforme, décomposition par pilier,
%     allocation (Shapley / Euler), priorisation de remédiation.
%  5. Bootstrap à deux niveaux : l'incertitude de calibration seule fait
%     un facteur ~2,5 sur la VaR. À afficher, c'est la limite honnête.
%  6. CADRAGE À TENIR : le NIVEAU absolu du SCR est illustratif (entité
%     notionnelle, sévérité issue de pertes US de grandes institutions).
%     Ce qui est vendu, c'est l'EFFET RELATIF et la hiérarchie des piliers.
